\section{Conclusion}\label{sec:conclusion}

In this study, we developed a computational framework that supplied evolving IV inputs for American-option valuation along separately generated stock paths. Fitted neural surfaces set contract-specific variance targets, and square-root factors translated those targets into lattice inputs. Greater surface flexibility and narrower parameter sharing improved calibration fit, while chronological evaluation exposed uneven transfer to later observations. The paired ablation showed that changing the IV process altered interim marks and liquidation outcomes on identical stock scenarios, with terminal payoffs remaining fixed. The later option forecasts did not establish a consistent advantage for the coupled factor over frozen IV. Supplying subsequently observed stock paths reduced option errors in a diagnostic that used future information, which motivated a separate examination of the stock forecasts. The larger 2025 stock comparison and the 2026 case study showed central errors close to an unchanged-price benchmark. Adaptive volatility improved LLY's distribution scores in both periods but not GS's, and earlier validation did not support a directional adjustment from the tested return features. These findings support the framework as a reproducible way to compare IV and stock assumptions, while limiting claims about their forecasting accuracy. The original JumpHMM remains the worked example rather than a newly revised stock estimator. Saved observations, fitted models, simulation outputs, and scoring code make the comparisons reviewable. Incomplete monthly contract histories, dependent evaluation dates, limited model selection, and residual strike inconsistencies remain constraints on broader use.
